% Original analytic teaching diagram, 2026-11-12.
% Uniform point in the unit disk: {R<=r} is the inner disk (area ratio r^2),
% {r<R<=r+h} is a thin ring of area pi((r+h)^2-r^2) ~ 2 pi r h.
\definecolor{VB}{HTML}{7188A3}\definecolor{VBF}{HTML}{E5EAF0}
\definecolor{VG}{HTML}{829C8B}\definecolor{VGF}{HTML}{E8EEE9}
\definecolor{VY}{HTML}{B4A16F}\definecolor{VYF}{HTML}{F1ECDD}
\definecolor{VR}{HTML}{A7787D}\definecolor{VRF}{HTML}{F0E5E6}
\definecolor{VP}{HTML}{9588A5}\definecolor{VPF}{HTML}{ECE8F0}
\definecolor{VI}{HTML}{3B4252}\definecolor{VM}{HTML}{6B7280}\definecolor{VL}{HTML}{4C4D4F}
\begin{tikzpicture}[x=1mm,y=1mm,>=Stealth,font=\figurefont\fontsize{8}{10}\selectfont,
lab/.style={text=VM,font=\figurefont\fontsize{7.5}{9}\selectfont}]
% panel (a): inner disk
\begin{scope}[shift={(22,24)}]
  \fill[VBF] (0,0) circle (12);
  \draw[VB,line width=1pt] (0,0) circle (12);
  \draw[VL,line width=.8pt] (0,0) circle (20);
  \fill[VL] (0,0) circle (.7);
  \draw[->,VL,line width=.6pt] (0,0) -- (-30:12) node[pos=.55,below,text=VI] {$r$};
  \draw[->,VL,line width=.6pt] (0,0) -- (60:20) node[pos=.8,right,text=VI] {$1$};
  \fill[VR] (110:7) circle (.9);
  \node[text=VI] at (0,-25) {(a)};
\end{scope}
% panel (b): thin ring
\begin{scope}[shift={(74,24)}]
  \fill[VBF,even odd rule] (0,0) circle (14) (0,0) circle (11);
  \draw[VB,line width=1pt] (0,0) circle (14);
  \draw[VB,line width=1pt] (0,0) circle (11);
  \draw[VL,line width=.8pt] (0,0) circle (20);
  \fill[VL] (0,0) circle (.7);
  \draw[->,VL,line width=.6pt] (0,0) -- (200:11) node[pos=.5,above,text=VI] {$r$};
  \draw[VR,line width=1.4pt] (-30:11) -- (-30:14);
  \node[text=VI] at (-30:16.6) {$h$};
  \fill[VR] (40:12.5) circle (.9);
  \node[text=VI] at (0,-25) {(b)};
\end{scope}
\end{tikzpicture}
